\appendix		%Latex sagen, dass ab hier nur noch Anhänge kommen
\chapter*{Anhänge}\thispagestyle{rmplain}	%erstelltes Kapitel welches aber nicht im Inhaltsverzeichnis aufscheint 
\addcontentsline{toc}{chapter}{Anhänge}		%Inhaltsverzeichniseintrag ohne nummerierung für zuvor erstelles Kapitel erzeugen
\label{anhaenge}


\renewcommand{\thesection}{\Alph{section}}					%Sections mit Großbuchstaben durchnummerieren
\renewcommand{\thefigure}{\Alph{section}.\arabic{figure}}	%Art der Bildbeschriftung auf Buchstaben.Nummerierung.: ändern


\section{Name von Anhang A}\thispagestyle{rmplain}
\label{app_a}
\setcounter{figure}{0}		%Muss bei jeder neuen Section gesetzt werden, damit die Caption wieder bei eins zu nummerieren beginnt (z.B. A.1. A.2. ... B.1. B.2.)

%Der Anhang kommt hier rein
\begin{figure}[htb]
	\centering  
	\includegraphics[scale=0.5]{img/Lord_of_The_Rings_Logo.png}
	\caption{Lord of The Rings Logo}
	\label{fig:lotrl1}
\end{figure}

%\begin{figure}[htb]
%	\centering  
%	\includegraphics[scale=0.5]{img/Lord_of_The_Rings_Logo.png}
%	\caption{Lord of The Rings Logo}
%	\label{fig:lotrl2}
%\end{figure}
%\newpage

%%\section{Name von Anhang B}\thispagestyle{rmplain}
%\label{app_b}
%\setcounter{figure}{0}		%Muss bei jeder neuen Section gesetzt werden, damit die Caption wieder bei eins zu nummerieren beginnt (z.B. A.1. A.2. ... B.1. B.2.)

%Der Anhang kommt hier rein
%\begin{figure}[htb]
%%	\includegraphics[scale=0.5]{img/Lord_of_The_Rings_Logo.png}
	%\caption{Lord of The Rings Logo}
	%\label{fig:lotrl3}
%\end{figure}
\newpage